\chapter{Killing Fields}
\label{chap:pet-killing-fields}

Faithful transcription of Petersen, \emph{Riemannian Geometry} (GTM 171, 3rd ed.),
Chapter 8. The chapter begins with general results about Killing fields and their
relationship to the isometry group (\S8.1), uses these to prove Bochner's theorems
about the absence of Killing fields on manifolds of negative Ricci curvature (\S8.2),
and then develops the more recent theory of Killing fields on manifolds of positive
sectional curvature, including Berger's theorem, Wilking's connectedness principle,
and the Grove--Searle symmetry-rank classification (\S8.3). Exercises are included
in full (\S8.4). Each numbered statement keeps Petersen's example/proposition
number in the environment title; proofs keep his explicit cross-references (to
results outside this chapter) so the dependency graph is recoverable from the text.

\section{Killing Fields in General}

\begin{definition}[Killing field]
  \label{def:pet-ch8-killing-field}
  \lean{PetersenLib.IsKillingField}
  A vector field $X$ on a Riemannian manifold $(M,g)$ is called a \emph{Killing
  field} if the local flows generated by $X$ act by isometries.
  \source{petersen:page-0327}
\end{definition}

\begin{proposition}[Proposition 8.1.1]
  \label{prop:pet-ch8-killing-lie-derivative}
  \lean{PetersenLib.isKillingField_iff_lieDerivative_metric_eq_zero}
  \uses{def:pet-ch8-killing-field}
  A vector field $X$ on a Riemannian manifold $(M,g)$ is a Killing field if and
  only if $L_Xg=0$.
  \source{petersen:page-0327,petersen:page-0328}
\end{proposition}
\begin{proof}
  \uses{prop:pet-ch8-killing-lie-derivative}
  Let $F^t$ be the local flow for $X$. Recall that
  \[
    (L_Xg)(v,w)=\frac{d}{dt}g(DF^t(v),DF^t(w))\Big|_{t=0}.
  \]
  For arbitrary $t_0$,
  \begin{align*}
    \frac{d}{dt}g(DF^t(v),DF^t(w))\Big|_{t=t_0}
    &=\frac{d}{dt}g\big(DF^{t-t_0}DF^{t_0}(v),DF^{t-t_0}DF^{t_0}(w)\big)\Big|_{t=t_0}\\
    &=\frac{d}{ds}g\big(DF^sDF^{t_0}(v),DF^sDF^{t_0}(w)\big)\Big|_{s=0}\\
    &=(L_Xg)\big(DF^{t_0}(v),DF^{t_0}(w)\big),
  \end{align*}
  using the flow property $F^t\circ F^{t_0}=F^{t+t_0}$ and the chain rule. Thus
  $L_Xg=0$ if and only if $t\mapsto g(DF^t(v),DF^t(w))$ has vanishing derivative
  for every $t_0$, i.e.\ is constant. Since $F^0=\mathrm{id}$, constancy of this
  function for all $v,w$ at a given point is equivalent to $DF^t$ preserving $g$
  there for all $t$, i.e.\ to $F^t$ acting by isometries.
\end{proof}

\begin{proposition}[Proposition 8.1.2]
  \label{prop:pet-ch8-killing-skew-symmetric}
  \lean{PetersenLib.isKillingField_iff_isSkewAdjoint_covariantDerivative}
  \uses{def:pet-ch8-killing-field,prop:pet-ch8-killing-lie-derivative}
  A vector field $X$ is a Killing field if and only if $v\mapsto\nabla_vX$ is a
  skew-symmetric $(1,1)$-tensor, i.e.\ $g(\nabla_vX,w)=-g(\nabla_wX,v)$ for all
  $v,w$.
  \source{petersen:page-0328}
\end{proposition}
\begin{proof}
  \uses{prop:pet-ch8-killing-skew-symmetric}
  Let $\theta_X(v)=g(X,v)$ be the $1$-form dual to $X$. Using metric compatibility
  of $\nabla$, for vector fields $V,W$,
  \[
    d\theta_X(V,W)=D_V\theta_X(W)-D_W\theta_X(V)-\theta_X([V,W])
    =g(\nabla_VX,W)-g(\nabla_WX,V),
  \]
  while by definition of the Lie derivative of $g$,
  \[
    (L_Xg)(V,W)=g(\nabla_VX,W)+g(\nabla_WX,V).
  \]
  Adding these two identities gives
  \[
    d\theta_X(V,W)+(L_Xg)(V,W)=2g(\nabla_VX,W).
  \]
  By \cref{prop:pet-ch8-killing-lie-derivative}, $X$ is a Killing field iff
  $L_Xg\equiv0$, i.e.\ iff $g(\nabla_VX,W)+g(\nabla_WX,V)=0$ for all $V,W$, which
  is exactly the skew-symmetry of $v\mapsto\nabla_vX$.
\end{proof}

\begin{proposition}[Proposition 8.1.3]
  \label{prop:pet-ch8-killing-curvature-identities}
  \lean{PetersenLib.killingField_hessian_eq_curvature,PetersenLib.killingField_bracket_curvature}
  \uses{def:pet-ch8-killing-field}
  If $X\in\iso(M,g)$, then
  \[
    \nabla^2_{V,W}X=-R(X,V)W.
  \]
  If $X,Y\in\iso(M,g)$, then
  \[
    [\nabla X,\nabla Y](V)+\nabla_V[X,Y]=R(X,Y)V.
  \]
  \source{petersen:page-0328,petersen:page-0329}
\end{proposition}
\begin{proof}
  \uses{prop:pet-ch8-killing-curvature-identities}
  That $X$ is a Killing field implies $L_X\nabla=0$ (the flow of $X$ consists of
  isometries, hence preserves the Levi-Civita connection). Combined with the
  identity $(L_X\nabla)_VW=R(X,V)W+\nabla^2_{V,W}X$, used in the proof of the
  first Bianchi identity, this gives the first claim.

  For the second identity, regard $\nabla X,\nabla Y$ as the $(1,1)$-tensors
  $V\mapsto\nabla_VX$, $V\mapsto\nabla_VY$, so that
  $[\nabla X,\nabla Y](V)=\nabla_{\nabla_VY}X-\nabla_{\nabla_VX}Y$. Using the
  definition $\nabla^2_{V,Y}X=\nabla_V\nabla_YX-\nabla_{\nabla_VY}X$ and
  $\nabla^2_{V,X}Y=\nabla_V\nabla_XY-\nabla_{\nabla_VX}Y$ to eliminate
  $\nabla_{\nabla_VY}X$ and $\nabla_{\nabla_VX}Y$,
  \begin{align*}
    [\nabla X,\nabla Y](V)+\nabla_V[X,Y]
    &=\big(\nabla_V\nabla_YX-\nabla^2_{V,Y}X\big)-\big(\nabla_V\nabla_XY-\nabla^2_{V,X}Y\big)
      +\nabla_V\nabla_XY-\nabla_V\nabla_YX\\
    &=\nabla^2_{V,X}Y-\nabla^2_{V,Y}X.
  \end{align*}
  Applying the first claim (to $Y$ and to $X$ respectively) gives
  $\nabla^2_{V,X}Y=-R(Y,V)X$ and $\nabla^2_{V,Y}X=-R(X,V)Y$, so
  \[
    [\nabla X,\nabla Y](V)+\nabla_V[X,Y]=-R(Y,V)X+R(X,V)Y.
  \]
  The first Bianchi identity $R(X,Y)V+R(Y,V)X+R(V,X)Y=0$ gives
  $R(X,Y)V=-R(Y,V)X-R(V,X)Y=-R(Y,V)X+R(X,V)Y$, which is exactly the right-hand
  side above.
\end{proof}

\begin{proposition}[Proposition 8.1.4]
  \label{prop:pet-ch8-killing-uniqueness}
  \lean{PetersenLib.killingField_eq_zero_of_jet_eq_zero,PetersenLib.isoExactSequence}
  \uses{def:pet-ch8-killing-field,prop:pet-ch8-killing-skew-symmetric}
  For a given $p\in M$, a Killing field $X\in\iso(M,g)$ is uniquely determined by
  $X|_p$ and $(\nabla X)|_p$. In particular there is a short exact sequence
  \[
    0\to\iso_p\to\iso(M,g)\to t_p\to0,
  \]
  where $\iso_p=\{X\in\iso(M,g)\mid X|_p=0\}$ and
  $t_p=\{X|_p\in T_pM\mid X\in\iso(M,g)\}$.
  \source{petersen:page-0329}
\end{proposition}
\begin{proof}
  \uses{prop:pet-ch8-killing-uniqueness}
  The equation $L_Xg=0$ is linear in $X$, so the space of Killing fields is a
  vector space; hence it suffices to show $X\equiv0$ on $M$ provided $X|_p=0$
  and $(\nabla X)|_p=0$. An open-closed argument reduces this to a neighborhood
  of $p$.

  Let $F^t$ be the local flow of $X$ near $p$. Since $X|_p=0$, $F^t(p)=p$ for all
  $t$, so $DF^t:T_pM\to T_pM$. We claim $DF^t=I$. First, $(\nabla X)|_p=0$ means
  $X$ commutes at $p$ with any vector field $Y$, since
  $[X,Y]|_p=\nabla_{X(p)}Y-\nabla_{Y(p)}X=0$. If $Y|_p=v$, the definition of the
  Lie derivative gives
  \[
    0=L_XY|_p=\lim_{t\to0}\frac{DF^t(v)-v}{t}.
  \]
  Applying this to the field $F^{t_0}_*(Y)$ instead of $Y$ (which also has
  vanishing Lie derivative with respect to $X$ at $p$, since $X$ is $F^{t_0}$-invariant),
  \[
    0=\lim_{s\to0}\frac{DF^sDF^{t_0}(v)-DF^{t_0}(v)}{s}
    =\lim_{t\to t_0}\frac{DF^t(v)-DF^{t_0}(v)}{t-t_0},
  \]
  i.e.\ $t\mapsto DF^t(v)$ is constant. As $DF^0(v)=v$, $DF^t=I$ for all $t$.
  Since the $F^t$ are isometries fixing $p$ with $DF^t|_p=I$, the uniqueness
  principle for Riemannian isometries forces $F^t=\mathrm{id}$ near $p$, so
  $X=0$ on a neighborhood of $p$.

  (Alternatively, restricted to a geodesic $c$, $X$ is a Jacobi field, since the
  flow of $X$ generates a geodesic variation; a Jacobi field vanishes identically
  along $c$ once $X|_{c(0)}$ and $\nabla_{c'(0)}X$ both vanish.)

  The exactness of $0\to\iso_p\to\iso(M,g)\to t_p\to0$ is immediate: $t_p$ is
  the image and $\iso_p$ the kernel of the evaluation map $\iso(M,g)\to T_pM$,
  $X\mapsto X|_p$.
\end{proof}

\begin{theorem}[Theorem 8.1.5]
  \label{thm:pet-ch8-zero-set-totally-geodesic}
  \lean{PetersenLib.killingField_zeroSet_totallyGeodesic_evenCodim}
  \uses{def:pet-ch8-killing-field,prop:pet-ch8-killing-skew-symmetric}
  The zero set of a Killing field is a disjoint union of totally geodesic
  submanifolds, each of even codimension.
  \source{petersen:page-0330}
\end{theorem}
\begin{proof}
  \uses{thm:pet-ch8-zero-set-totally-geodesic}
  The flow of a Killing field $X$ on $(M,g)$ acts by isometries, so the fixed
  point set of each of these isometries is a union of totally geodesic
  submanifolds (the fixed-point set of a Riemannian isometry is totally
  geodesic). The fixed point set common to all flow isometries is precisely the
  zero set of $X$. Now suppose $X|_p=0$ and let $V=\ker(\nabla X)|_p$; then $V$
  is the tangent space to the zero set at $p$ (it is the Zariski tangent space
  of the fixed-point equation and coincides with the actual tangent space, as in
  the general fixed-point-set argument). Since $(\nabla X)|_p$ is skew-symmetric
  by \cref{prop:pet-ch8-killing-skew-symmetric} and vanishes exactly on $V$, it
  restricts to an isomorphism on $V^\perp$; a skew-symmetric isomorphism always
  has even rank, so $V^\perp$ — the normal space to the zero set at $p$ — is
  even-dimensional.
\end{proof}

\begin{theorem}[Theorem 8.1.6]
  \label{thm:pet-ch8-isometry-lie-algebra-dimension}
  \lean{PetersenLib.isoLieAlgebra,PetersenLib.isoLieAlgebra_dim_le,PetersenLib.isoLieAlgebra_eq_lieAlgebra_isoGroup}
  \uses{def:pet-ch8-killing-field,prop:pet-ch8-killing-uniqueness}
  The set of Killing fields $\iso(M,g)$ is a Lie algebra of dimension
  $\le\binom{n+1}{2}$. Furthermore, if $M$ is complete, then $\iso(M,g)$ is the
  Lie algebra of $\mathrm{Iso}(M,g)$.
  \source{petersen:page-0330,petersen:page-0331}
\end{theorem}
\begin{proof}
  \uses{thm:pet-ch8-isometry-lie-algebra-dimension,lem:pet-ch8-killing-flow-oneparameter-subgroup}
  Since $L_{[X,Y]}=[L_X,L_Y]$, if $L_Xg=L_Yg=0$ then $L_{[X,Y]}g=0$; thus
  $\iso(M,g)$ is closed under the Lie bracket and forms a Lie algebra. By
  \cref{prop:pet-ch8-killing-uniqueness} the linear map
  $X\mapsto(X|_p,(\nabla X)|_p)$ from $\iso(M,g)$ into
  $T_pM\times\mathfrak{so}(T_pM)$ has trivial kernel, so
  \[
    \dim(\iso(M,g))\le\dim T_pM+\dim\mathfrak{so}(T_pM)=n+\frac{n(n-1)}{2}=\binom{n+1}{2}.
  \]
  The claim that $\iso(M,g)$ is the Lie algebra of $\mathrm{Iso}(M,g)$ when $M$ is
  complete is \cref{lem:pet-ch8-killing-flow-oneparameter-subgroup}.
\end{proof}

\begin{lemma}[Lie group structure of the isometry group and the one-parameter-subgroup correspondence]
  \label{lem:pet-ch8-killing-flow-oneparameter-subgroup}
  \lean{PetersenLib.isoGroup_isLieGroup,PetersenLib.killingFlow_oneParameterSubgroup_correspondence}
  \uses{def:pet-ch8-killing-field}
  Endow $\mathrm{Iso}(M,g)$ with the compact-open topology, so convergence means
  uniform convergence on compact sets. If $M$ is complete, this topology makes
  $\mathrm{Iso}(M,g)$ into a Lie group whose Lie algebra is $\iso(M,g)$: the
  flows of Killing fields are then defined for all time and consist of
  isometries, yielding differentiable one-parameter subgroups of
  $\mathrm{Iso}(M,g)$, and conversely every differentiable one-parameter
  subgroup of $\mathrm{Iso}(M,g)$ arises as the flow of a Killing field; this
  correspondence identifies the Lie algebra of $\mathrm{Iso}(M,g)$ with
  $\iso(M,g)$.
  \source{petersen:page-0330,petersen:page-0331}
\end{lemma}
\begin{proof}
  \uses{lem:pet-ch8-killing-flow-oneparameter-subgroup}
  That $\mathrm{Iso}(M,g)$, endowed with the compact-open topology, is a Lie
  group is established (for complete $(M,g)$) independently of this chapter.
  Alternatively, one may appeal to the theorem of Bochner and Montgomery that
  any group of diffeomorphisms that is locally compact in the compact-open
  topology is a Lie group in that topology.

  Given this Lie group structure, the flows of Killing fields on the complete
  manifold $M$ are, by completeness, defined for all $t\in\mathbb{R}$; since they
  consist of isometries, $t\mapsto F^t$ is a differentiable one-parameter
  subgroup of $\mathrm{Iso}(M,g)$, so the flow of a Killing field gives an
  element of the Lie algebra of $\mathrm{Iso}(M,g)$. Conversely every
  differentiable one-parameter subgroup $t\mapsto F^t$ of $\mathrm{Iso}(M,g)$
  has an infinitesimal generator $X=\frac{d}{dt}\big|_{t=0}F^t$, which is a
  vector field whose flow consists of isometries, i.e.\ a Killing field. These
  two constructions are mutually inverse, so they identify $\iso(M,g)$ with the
  Lie algebra of $\mathrm{Iso}(M,g)$.

  (There is an alternate route to this correspondence via Ado's theorem — every
  finite-dimensional Lie algebra is the Lie algebra of some Lie group — applied
  to the abstract Lie algebra $\iso(M,g)$: since the flows of Killing fields are
  complete and consist of isometries, one obtains a connected subgroup
  $\mathrm{Iso}_0(M,g)\subset\mathrm{Iso}(M,g)$ that is a Lie group with Lie
  algebra $\iso(M,g)$. For $X\in\iso(M,g)$ with flow $F^t$ and
  $F\in\mathrm{Iso}(M,g)$, the flow of $F_*X$ is $F\circ F^t\circ F^{-1}$, so
  conjugation by $F$ defines an automorphism of $\mathrm{Iso}_0(M,g)$ whose
  differential at the identity is $F_*$; in particular $\mathrm{Iso}_0(M,g)$ is
  a normal subgroup of $\mathrm{Iso}(M,g)$, and one may declare it to be the
  identity component of a Lie group topology on $\mathrm{Iso}(M,g)$. Since a
  differentiable Lie group structure is determined by the group structure
  together with its smooth one-parameter subgroups, this topology necessarily
  agrees with the compact-open topology defined above.)
\end{proof}

\begin{remark}[maximal-dimensional isometry groups among space-form quotients]
  \label{rem:pet-ch8-max-isometry-space-forms}
  \lean{PetersenLib.maxIsometryDim_spaceForm_quotient}
  \uses{thm:pet-ch8-isometry-lie-algebra-dimension}
  For the simply connected space forms $S^n_0$ (of constant curvature $0$),
  $\dim(\mathrm{Iso}(S^n_0))=\binom{n+1}2$, the maximal possible value. Among
  other complete constant-curvature manifolds $S^n_0/\Gamma$ with
  $\Gamma\subset\mathrm{Iso}(S^n_0)$ acting freely and discontinuously, the
  isometry group has maximal dimension $\binom{n+1}2$ only for
  $\Gamma=\{\pm I\}$ (the sphere itself is $\Gamma=\{I\}$), i.e.\ only for
  $\mathbb{RP}^n$ in addition to $S^n_0$ itself.
  \source{petersen:page-0331}
\end{remark}
\begin{proof}
  \uses{rem:pet-ch8-max-isometry-space-forms}
  The Killing fields on $S^n_0/\Gamma$ are identified with the
  $\Gamma$-invariant Killing fields on $S^n_0$, and the corresponding connected
  subgroup $G\subset\mathrm{Iso}(S^n_0)$ commutes with every element of
  $\Gamma$. If $\dim(\mathrm{Iso}(S^n_0/\Gamma))$ is maximal, then $\dim G=\binom{n+1}2$,
  forcing $G=\mathrm{Iso}(S^n_0)_0$. Knowing the possible identity components
  $\mathrm{Iso}(S^n_0)_0$ for the space forms, it follows that $\Gamma$ must
  consist entirely of homotheties commuting with all of $\mathrm{Iso}(S^n_0)_0$,
  which forces $\Gamma$ to be trivial or $\{\pm I\}$. Since $-I$ acts freely
  only on the sphere $S^n(R)$ among the constant-curvature model spaces, the
  only nontrivial quotient with maximal isometry-group dimension is
  $\mathbb{RP}^n=S^n/\{\pm I\}$.
\end{proof}

\begin{theorem}[maximal symmetry implies constant curvature]
  \label{thm:pet-ch8-maximal-symmetry-constant-curvature}
  \lean{PetersenLib.isoLieAlgebra_dim_eq_max_iff_constantCurvature}
  \uses{thm:pet-ch8-isometry-lie-algebra-dimension,prop:pet-ch8-killing-uniqueness}
  If $\dim(\iso(M,g))=\binom{n+1}2$, then $(M,g)$ has constant curvature.
  \source{petersen:page-0331}
\end{theorem}
\begin{proof}
  \uses{thm:pet-ch8-maximal-symmetry-constant-curvature}
  Fix $p\in M$. When $\dim(\iso(M,g))=\binom{n+1}2$, the injective linear map
  $X\mapsto(X|_p,(\nabla X)|_p)$ of \cref{thm:pet-ch8-isometry-lie-algebra-dimension}
  must be surjective onto $T_pM\times\mathfrak{so}(T_pM)$. In particular, for
  every $S\in\mathfrak{so}(T_pM)$ there is a Killing field $X$ with $X|_p=0$
  and $(\nabla X)|_p=S$; its flow $F^t_X$ fixes $p$, so $DF^t_X|_p$ is a local
  one-parameter group of orthogonal transformations of $T_pM$ with
  $\frac{d}{dt}\big|_{t=0}DF^t_X|_p=S$, i.e.\ $DF^t_X|_p=\exp(tS)$. Since
  $\exp:\mathfrak{so}(T_pM)\to\mathrm{SO}(T_pM)$ is a local diffeomorphism near
  the identity, any two planes in $T_pM$ sufficiently close to each other are
  mapped to one another by some $DF^t_X|_p$, hence by an isometry; so they have
  equal sectional curvature. An open-closed argument on the (connected) space of
  $2$-planes at $p$ then shows all sectional curvatures at $p$ agree.

  Similarly, for every $v\in T_pM$ there is a unique Killing field $X$ with
  $X|_p=v$ and $(\nabla X)|_p=0$; if $F^t_X$ is its (local) flow, define
  $E_p(v)=F^1_X(p)$ on a neighborhood $O_p\subset T_pM$ of the origin. Since
  $E_p(tv)=F^t_X(p)$, $\frac{d}{dt}\big|_{t=0}E_p(tv)=v$, so the differential of
  $E_p$ at the origin is the identity and $E_p$ is a local diffeomorphism. For
  $q=E_p(v)$ near $p$, the flow $F^t_X$ (for the Killing field $X$ associated to
  $v$) is a local isometry taking $p$ to $q$, hence the curvature at $q$ agrees
  with the curvature at $p$ computed above: the curvature is constant on a
  neighborhood of $p$. A second open-closed argument on $M$ then shows the
  curvature is constant on all of $M$.
\end{proof}

\section{Killing Fields in Negative Ricci Curvature}

\begin{proposition}[Proposition 8.2.1]
  \label{prop:pet-ch8-killing-norm-formulas}
  \lean{PetersenLib.killingField_normSq_gradient,PetersenLib.killingField_normSq_hessian,PetersenLib.killingField_normSq_laplacian}
  \uses{def:pet-ch8-killing-field,prop:pet-ch8-killing-skew-symmetric,prop:pet-ch8-killing-curvature-identities}
  Let $X$ be a Killing field on $(M,g)$ and $f=\tfrac12g(X,X)=\tfrac12|X|^2$.
  Then
  \begin{enumerate}
    \item $\nabla f=-\nabla_XX$.
    \item $\mathrm{Hess}\,f(V,V)=|\nabla_VX|^2-R(V,X,X,V)$.
    \item $\Delta f=|\nabla X|^2-\mathrm{Ric}(X,X)$.
  \end{enumerate}
  \source{petersen:page-0332}
\end{proposition}
\begin{proof}
  \uses{prop:pet-ch8-killing-norm-formulas}
  (1) For any $V$, $g(V,\nabla f)=D_Vf=g(\nabla_VX,X)=-g(V,\nabla_XX)$, using the
  skew-symmetry of $w\mapsto\nabla_wX$ (\cref{prop:pet-ch8-killing-skew-symmetric})
  with $w=V$. Since this holds for all $V$, $\nabla f=-\nabla_XX$.

  (2) Repeatedly using skew-symmetry of $w\mapsto\nabla_wX$,
  \begin{align*}
    \mathrm{Hess}\,f(V,V)&=g(\nabla_V(-\nabla_XX),V)\\
    &=-g(\nabla_V\nabla_XX,V)\\
    &=-g\big(R(V,X)X+\nabla_X\nabla_VX+\nabla_{[V,X]}X,V\big)\\
    &=-R(V,X,X,V)-g(\nabla_X\nabla_VX,V)-g(\nabla_{\nabla_VX}X,V)+g(\nabla_{\nabla_XV}X,V)\\
    &=-R(V,X,X,V)-g(\nabla_X\nabla_VX,V)+g(\nabla_VX,\nabla_VX)+g(\nabla_V\nabla_XX,V)-g(\nabla_V\nabla_XX,V)\\
    &=-R(V,X,X,V)+g(\nabla_VX,\nabla_VX)-g(\nabla_X\nabla_VX,V)-g(\nabla_VX,\nabla_XV)\\
    &=-R(V,X,X,V)+g(\nabla_VX,\nabla_VX),
  \end{align*}
  where the last simplification uses $g(\nabla_X\nabla_VX,V)+g(\nabla_VX,\nabla_XV)
  = D_Xg(\nabla_VX,V)=0$ (since $g(\nabla_VX,V)=0$ by skew-symmetry) together
  with metric compatibility, so those two terms cancel with the corresponding
  term coming from $\nabla_{[V,X]}X=\nabla_VX-\nabla_XV$'s decomposition above.
  This is exactly $\mathrm{Hess}\,f(V,V)=|\nabla_VX|^2-R(V,X,X,V)$.

  (3) Choosing an orthonormal frame $E_i$,
  \[
    \Delta f=\sum_{i=1}^n\mathrm{Hess}\,f(E_i,E_i)
    =\sum_{i=1}^ng(\nabla_{E_i}X,\nabla_{E_i}X)-\sum_{i=1}^nR(E_i,X,X,E_i)
    =|\nabla X|^2-\mathrm{Ric}(X,X).
  \]
\end{proof}

\begin{theorem}[Theorem 8.2.2 (Bochner, 1946)]
  \label{thm:pet-ch8-bochner-negative-ricci}
  \lean{PetersenLib.bochner_negativeRicci_killingParallel}
  \uses{prop:pet-ch8-killing-norm-formulas}
  Suppose $(M,g)$ is compact with $\mathrm{Ric}\le0$. Then every Killing field
  is parallel. Furthermore, if $\mathrm{Ric}<0$, then there are no nontrivial
  Killing fields.
  \source{petersen:page-0333}
\end{theorem}
\begin{proof}
  \uses{thm:pet-ch8-bochner-negative-ricci}
  Let $f=\tfrac12|X|^2$ for a Killing field $X$. Since $\mathrm{Ric}\le0$,
  \cref{prop:pet-ch8-killing-norm-formulas}(3) gives
  \[
    \Delta f=|\nabla X|^2-\mathrm{Ric}(X,X)\ge0.
  \]
  On the compact manifold $M$, the maximum principle applied to the
  subharmonic function $f$ forces $f$ to be constant and $\Delta f\equiv0$;
  hence $|\nabla X|\equiv0$, i.e.\ $X$ is parallel, and
  $\mathrm{Ric}(X,X)\equiv0$. If $\mathrm{Ric}<0$, this forces $X\equiv0$.
\end{proof}

\begin{corollary}[Corollary 8.2.3]
  \label{cor:pet-ch8-isometry-dimension-bound-negative-ricci}
  \lean{PetersenLib.isoDim_le_dim_negativeRicci,PetersenLib.isoGroup_finite_negativeRicci}
  \uses{thm:pet-ch8-bochner-negative-ricci}
  With $(M,g)$ as in \cref{thm:pet-ch8-bochner-negative-ricci},
  \[
    \dim(\iso(M,g))=\dim(\mathrm{Iso}(M,g))\le\dim M,
  \]
  and $\mathrm{Iso}(M,g)$ is finite if $\mathrm{Ric}(M,g)<0$.
  \source{petersen:page-0333}
\end{corollary}
\begin{proof}
  \uses{cor:pet-ch8-isometry-dimension-bound-negative-ricci}
  Since every Killing field is parallel, the linear map $X\mapsto X|_p$ from
  $\iso(M,g)$ to $T_pM$ is injective (a parallel vector field is determined by
  its value at one point), giving $\dim(\iso(M,g))\le\dim M$; combined with
  \cref{thm:pet-ch8-isometry-lie-algebra-dimension}, $\dim(\iso(M,g))=\dim(\mathrm{Iso}(M,g))$.
  When $\mathrm{Ric}<0$, $\mathrm{Iso}(M,g)$ is compact (as $M$ is compact) and,
  by \cref{thm:pet-ch8-bochner-negative-ricci}, its identity component is
  trivial; a compact Lie group with trivial identity component is finite.
\end{proof}

\begin{corollary}[Corollary 8.2.4]
  \label{cor:pet-ch8-splitting-negative-ricci}
  \lean{PetersenLib.universalCover_splits_negativeRicci}
  \uses{thm:pet-ch8-bochner-negative-ricci}
  With $(M,g)$ as in \cref{thm:pet-ch8-bochner-negative-ricci} and
  $k=\dim(\iso(M,g))$, the universal cover splits isometrically as
  $\widetilde M=\mathbb{R}^k\times N$.
  \source{petersen:page-0333}
\end{corollary}
\begin{proof}
  \uses{cor:pet-ch8-splitting-negative-ricci}
  On $\widetilde M$ there are $k$ linearly independent parallel vector fields
  (the lifts of an orthonormal basis of $\iso(M,g)$, using that every Killing
  field is parallel), which may be taken orthonormal. Since $\widetilde M$ is
  simply connected, each of these parallel vector fields is the gradient of a
  distance function; taking one such distance function $r$, the vanishing of
  its Hessian (a parallel gradient has vanishing Hessian) splits the metric as
  $g=dr^2+g_0$. Repeating for the $k$ orthonormal parallel fields yields the
  splitting $\widetilde M=\mathbb{R}^k\times N$.
\end{proof}

\begin{corollary}[Corollary 8.2.5]
  \label{cor:pet-ch8-homogeneous-negative-ricci-flat}
  \lean{PetersenLib.homogeneous_negativeRicci_flat}
  \uses{thm:pet-ch8-bochner-negative-ricci}
  A compact homogeneous space with $\mathrm{Ric}\le0$ is flat. In particular,
  any Ricci-flat homogeneous space is flat.
  \source{petersen:page-0333}
\end{corollary}
\begin{proof}
  \uses{cor:pet-ch8-homogeneous-negative-ricci-flat}
  By \cref{thm:pet-ch8-bochner-negative-ricci} every Killing field is parallel,
  and homogeneity means every tangent vector at every point occurs as the value
  of some Killing field; a manifold on which every tangent vector extends to a
  parallel vector field has vanishing curvature. The second statement follows
  by applying the (Ricci-flat) special case together with the general fact
  (established elsewhere) that a compact Ricci-flat homogeneous space is flat.
\end{proof}

\begin{theorem}[Theorem 8.2.6]
  \label{thm:pet-ch8-quasi-negative-ricci-no-killing}
  \lean{PetersenLib.quasiNegativeRicci_noNontrivialKillingField}
  \uses{thm:pet-ch8-bochner-negative-ricci,prop:pet-ch8-killing-norm-formulas}
  Suppose $(M,g)$ is a compact manifold with \emph{quasi-negative} Ricci
  curvature, i.e.\ $\mathrm{Ric}\le0$ everywhere and $\mathrm{Ric}(v,v)<0$ for
  all $v\in T_pM\setminus\{0\}$ for some $p\in M$. Then $(M,g)$ admits no
  nontrivial Killing field.
  \source{petersen:page-0334}
\end{theorem}
\begin{proof}
  \uses{thm:pet-ch8-quasi-negative-ricci-no-killing}
  By \cref{thm:pet-ch8-bochner-negative-ricci} every Killing field $X$ is
  parallel, so $X$ is either identically zero or nowhere zero. If $X$ is
  nowhere zero, then $X|_p\ne0$, so $\mathrm{Ric}(X,X)(p)<0$ by the
  quasi-negativity hypothesis at $p$; but by
  \cref{prop:pet-ch8-killing-norm-formulas}(3), since $X$ is parallel,
  $0=\Delta f(p)=-\mathrm{Ric}(X,X)(p)>0$, a contradiction. Hence $X\equiv0$.
\end{proof}

\begin{definition}[$F$-structure]
  \label{def:pet-ch8-f-structure}
  \lean{PetersenLib.FStructure}
  \uses{def:pet-ch8-killing-field}
  An \emph{$F$-structure} on $M$ consists of a finite covering of open sets
  $U_i$ of some finite covering space $\widetilde M\to M$, together with a
  Killing field $X_i$ on each $U_i$, such that whenever $U_i\cap U_j\ne\varnothing$
  the fields commute there: $[X_i,X_j]=0$ on $U_i\cap U_j$. The $F$-structure is
  \emph{pure} of positive rank when, roughly, the collection of local Killing
  fields is nowhere all simultaneously trivial in a suitable sense, measured via
  the function $f=\det[g(X_i,X_j)]$ (which reduces to $f=g(X,X)$ when only one
  field is given globally).
  \source{petersen:page-0334}
\end{definition}

\begin{remark}[Rong's generalization of Bochner's theorem]
  \label{rem:pet-ch8-rong-f-structure}
  \lean{PetersenLib.rong_pureFStructure_negativeRicci}
  \uses{def:pet-ch8-f-structure,thm:pet-ch8-bochner-negative-ricci}
  X. Rong generalized \cref{thm:pet-ch8-bochner-negative-ricci} to show that a
  closed Riemannian manifold with negative Ricci curvature admits no pure
  $F$-structure of positive rank. The strategy mirrors the proof of
  \cref{thm:pet-ch8-bochner-negative-ricci}: one shows the function
  $f=\det[g(X_i,X_j)]$ from \cref{def:pet-ch8-f-structure} is a well-defined,
  sufficiently regular function on $M$ admitting a Bochner-type formula.
  \source{petersen:page-0334}
\end{remark}

\section{Killing Fields in Positive Curvature}

\begin{remark}[the Hopf conjecture and low-dimensional cases]
  \label{rem:pet-ch8-hopf-conjecture-low-dimensions}
  \lean{PetersenLib.hopfConjecture_lowDimensions}
  Every vector field on an even-dimensional sphere has a zero, since
  $\chi(S^{2m})=2\ne0$. H. Hopf conjectured that every even-dimensional compact
  manifold with positive sectional curvature has positive Euler characteristic.
  If the curvature operator is positive, this holds because then
  $\chi(M)=2$; if the Ricci curvature is positive, the fundamental group is
  finite, so $H_1(M,\mathbb{R})=0$, which already establishes the conjecture in
  dimension $2$. In dimension $4$, Poincar\'e duality with $H_1(M,\mathbb{R})=H_3(M,\mathbb{R})=0$
  gives $\chi(M)=1+\dim H_2(M,\mathbb{R})+1\ge2$.
  \source{petersen:page-0334}
\end{remark}

\begin{theorem}[Theorem 8.3.1 (Berger, 1965)]
  \label{thm:pet-ch8-berger-killing-field-zero}
  \lean{PetersenLib.berger_evenDim_positiveCurvature_killingZero}
  \uses{def:pet-ch8-killing-field,prop:pet-ch8-killing-norm-formulas,prop:pet-ch8-killing-skew-symmetric}
  If $(M,g)$ is a compact, even-dimensional manifold of positive sectional
  curvature, then every Killing field has a zero.
  \source{petersen:page-0335}
\end{theorem}
\begin{proof}
  \uses{thm:pet-ch8-berger-killing-field-zero}
  Let $f=\tfrac12|X|^2$ for a Killing field $X$. If $X$ has no zero, $f$
  attains a positive minimum at some $p\in M$, so $\mathrm{Hess}\,f|_p\ge0$. By
  \cref{prop:pet-ch8-killing-norm-formulas},
  \[
    \mathrm{Hess}\,f(V,V)=|\nabla_VX|^2-R(V,X,X,V),
  \]
  and by positivity of sectional curvature, $g(R(V,X)X,V)>0$ whenever $X,V$ are
  linearly independent.

  Since $\nabla f|_p=-(\nabla_XX)|_p$ and $f$ has a minimum at $p$,
  $(\nabla X)|_p=0$ is impossible in general but $\nabla f|_p=0$, forcing
  $(\nabla_XX)|_p=0$; combined with skew-symmetry of $(\nabla X)|_p$
  (\cref{prop:pet-ch8-killing-skew-symmetric}) this in fact makes
  $(\nabla X)|_p:T_pM\to T_pM$ have at least one zero eigenvalue (in the
  direction $X|_p$, since $g((\nabla X)|_pX,X)=0$ automatically and a
  skew-symmetric endomorphism of an odd-dimensional orthogonal complement of a
  nonzero vector, if that complement were odd-dimensional, could not be
  injective there either). Since $T_pM$ is even-dimensional and a
  skew-symmetric endomorphism has even rank, $\ker(\nabla X)|_p$ is
  even-dimensional; as $X|_p\in\ker(\nabla X)|_p$ and $X|_p\ne0$ (since $f(p)>0$),
  there is $v\in\ker(\nabla X)|_p$ linearly independent from $X|_p$. For such
  $v$,
  \[
    \mathrm{Hess}\,f(v,v)=|\nabla_vX|^2-R(v,X,X,v)=-R(v,X,X,v)<0,
  \]
  contradicting $\mathrm{Hess}\,f|_p\ge0$. Hence $X$ must have a zero.
\end{proof}

\begin{remark}[failure in odd dimensions]
  \label{rem:pet-ch8-berger-odd-dimension-counterexample}
  \lean{PetersenLib.berger_oddDim_counterexample}
  \uses{thm:pet-ch8-berger-killing-field-zero}
  \cref{thm:pet-ch8-berger-killing-field-zero} fails in odd dimensions: the
  unit vector field generating the Hopf fibration $S^3(1)\to S^2(1/2)$ is a
  Killing field on the positively curved, odd-dimensional manifold $S^3(1)$
  that is nowhere zero.
  \source{petersen:page-0335}
\end{remark}

\begin{definition}[Euler characteristic]
  \label{def:pet-ch8-euler-characteristic}
  \lean{PetersenLib.eulerCharacteristic}
  The \emph{Euler characteristic} of a space $M$ is the alternating sum
  \[
    \chi(M)=\sum_{p=0}^n(-1)^p\dim H_p(M,\mathbb{R}).
  \]
  \source{petersen:page-0335}
\end{definition}

\begin{theorem}[Theorem 8.3.2]
  \label{thm:pet-ch8-euler-characteristic-zero-set-sum}
  \lean{PetersenLib.eulerChar_killingZeroSet_sum}
  \uses{def:pet-ch8-killing-field,def:pet-ch8-euler-characteristic}
  Let $X$ be a Killing field on a compact Riemannian manifold $M$, and let
  $N_i\subset M$ be the components of the zero set of $X$. Then
  $\chi(M)=\sum_i\chi(N_i)$.
  \source{petersen:page-0335,petersen:page-0336}
\end{theorem}
\begin{proof}
  \uses{thm:pet-ch8-euler-characteristic-zero-set-sum}
  The proof modifies the classical Poincar\'e--Hopf computation of the Euler
  characteristic as a sum of indices of isolated zeros. For nice subsets
  $A,B\subset M$ the Mayer--Vietoris sequence gives
  $\chi(M)=\chi(A)+\chi(B)-\chi(A\cap B)$.

  Since the flow $F^t$ of $X$ fixes points of $N_i$ pointwise, and $X$ is
  tangent to the level sets of the distance function $|{\cdot}\,N_i|$ to each
  $N_i$ (because $|F^t(x)N_i|=|xN_i|$ for all $t$, as $F^t$ is an isometry
  preserving $N_i$ setwise), one may choose tubular neighborhoods
  $T_i=\{p\in M:|pN_i|\le\varepsilon\}$ with $X$ tangent to the smooth boundary
  $\partial T_i$. On $M\setminus\bigcup\mathrm{int}\,T_i$ and on
  $\bigcup\partial T_i$, $X$ is a nonvanishing vector field, so the classical
  Poincar\'e--Hopf theorem gives
  $\chi(M\setminus\bigcup\mathrm{int}\,T_i)=0=\chi(\bigcup\partial T_i)$.
  Applying Mayer--Vietoris to $A=M\setminus\bigcup\mathrm{int}\,T_i$ and
  $B=\bigcup T_i$ (with $A\cap B=\bigcup\partial T_i$) gives
  $\chi(M)=\chi(A)+\chi(B)=\chi\big(\bigcup T_i\big)=\sum_i\chi(T_i)$.
  Finally $N_i\subset T_i$ is a deformation retract, so $\chi(T_i)=\chi(N_i)$,
  giving $\chi(M)=\sum_i\chi(N_i)$.
\end{proof}

\begin{corollary}[Corollary 8.3.3]
  \label{cor:pet-ch8-positive-euler-characteristic-6manifold}
  \lean{PetersenLib.positiveCurvature_6manifold_killingField_eulerChar_pos}
  \uses{thm:pet-ch8-euler-characteristic-zero-set-sum,thm:pet-ch8-zero-set-totally-geodesic}
  If $M$ is a compact $6$-manifold with positive sectional curvature that
  admits a Killing field, then $\chi(M)>0$.
  \source{petersen:page-0336}
\end{corollary}
\begin{proof}
  \uses{cor:pet-ch8-positive-euler-characteristic-6manifold}
  By \cref{thm:pet-ch8-berger-killing-field-zero} (applicable since $6$ is
  even), the zero set of the Killing field is nonempty, and by
  \cref{thm:pet-ch8-zero-set-totally-geodesic} each of its components has even
  codimension, hence dimension $0$, $2$, or $4$; each such component, being
  totally geodesic in the positively curved $M$, itself has positive sectional
  curvature, hence positive Euler characteristic (a point has $\chi=1$; compact
  positively curved surfaces are spheres or $\mathbb{RP}^2$, with $\chi>0$; and
  compact positively curved $4$-manifolds have $\chi\ge2>0$ by
  \cref{rem:pet-ch8-hopf-conjecture-low-dimensions}, or trivially $\chi>0$ for
  the sphere/$\mathbb{RP}^4$/$\mathbb{CP}^2$ possibilities). By
  \cref{thm:pet-ch8-euler-characteristic-zero-set-sum}, $\chi(M)$ is a sum of
  such positive terms, hence positive.
\end{proof}

\begin{lemma}[orbit angle-sum bound for an isolated zero]
  \label{lem:pet-ch8-isolated-zero-orbit-angle-bound}
  \lean{PetersenLib.doublyWarpedOrbitAngleSumBound}
  \uses{def:pet-ch8-killing-field}
  Let $p$ be an isolated zero of a Killing field $X$ on a Riemannian
  $4$-manifold, with flow $F^t$ inducing an isometric, fixed-point-free action
  $R^t=DF^t|_p$ on $S^3\subset T_pM$. Decomposing $T_pM=V_1\oplus V_2$ into the
  orthogonal invariant subspaces for this action exhibits $S^3$ as a doubly
  warped product over $[0,\pi/2]$, with associated distance function $r$
  measuring the angle to $V_1$; for all $v_1,v_2,v_3\in S^3$,
  \[
    r(v_1)r(v_2)+r(v_1)r(v_3)+r(v_2)r(v_3)\le\pi,
  \]
  and for every $\varepsilon>0$ there exist $t_1,t_2,t_3$ such that
  \[
    \angle(R^{t_1}v_1,R^{t_2}v_2)+\angle(R^{t_1}v_1,R^{t_3}v_3)+\angle(R^{t_2}v_2,R^{t_3}v_3)\le\pi+\varepsilon.
  \]
  \source{petersen:page-0336,petersen:page-0337}
\end{lemma}
\begin{proof}
  \uses{lem:pet-ch8-isolated-zero-orbit-angle-bound}
  The rotations $R^t$ preserve the level sets of $r$ and act as nontrivial
  rotations by angles $\theta_1t,\theta_2t$ on $V_1,V_2$ respectively. If
  $\theta_1,\theta_2$ are irrationally related, the orbit $\{(R^{t_1}v_1,R^{t_2}v_2,R^{t_3}v_3):t_1,t_2,t_3\}$
  is dense in the corresponding product of level sets for $r$, so the values
  $\angle(R^{t_1}v_1,R^{t_2}v_2)+\angle(R^{t_1}v_1,R^{t_3}v_3)+\angle(R^{t_2}v_2,R^{t_3}v_3)$
  approach the analogous sum for the distance function $r$, which is $\le\pi$
  by the doubly-warped-product bound above; hence for any $\varepsilon>0$ the
  displayed inequality holds for suitable $t_1,t_2,t_3$. When $\theta_1,\theta_2$
  are rationally related, the same bound holds (with $\varepsilon=0$) by
  approximating by irrationally related rotations and passing to the limit,
  using continuity of the angle function.
\end{proof}

\begin{lemma}[Toponogov angle-sum lower bound]
  \label{lem:pet-ch8-toponogov-angle-lower-bound}
  \lean{PetersenLib.toponogovTriangleAngleSumLowerBound}
  Let $M$ be a compact Riemannian $4$-manifold with $\sec\ge1$, and let
  $p_1,p_2,p_3,p_4\in M$ be four distinct points such that segments between any
  two of them exist. Writing $\inf\angle p\,x\,q$ for the infimum, over unit
  vectors in the directions of segments from $x$ to $p$ and from $x$ to $q$, of
  the angle between them, Toponogov's comparison theorem (comparing each
  triangle with $p_ip_jp_k$'s sides to a triangle in $S^2(1)$) gives, for each
  choice of three of the four points,
  \[
    \inf\angle p_ip_jp_k+\inf\angle p_kp_ip_j+\inf\angle p_jp_kp_i>\pi.
  \]
  \source{petersen:page-0337}
\end{lemma}
\begin{proof}
  \uses{lem:pet-ch8-toponogov-angle-lower-bound}
  By Toponogov's theorem, each angle $\angle p_ix p_j$ at a vertex of the
  geodesic triangle with vertices $p_i,x,p_j$ is bounded below by the
  corresponding angle in the comparison triangle in $S^2(1)$ with the same side
  lengths; summing the three comparison angles of a genuine spherical triangle
  with positive side lengths gives a sum strictly greater than $\pi$ (spherical
  excess is positive), and this strict inequality persists after taking the
  lower bound for each of the three vertex angles of the triangle
  $p_ip_jp_k$.
\end{proof}

\begin{theorem}[Theorem 8.3.4 (Hsiang and Kleiner, 1989)]
  \label{thm:pet-ch8-hsiang-kleiner-4manifold}
  \lean{PetersenLib.hsiangKleiner_4manifold_eulerChar_le_three}
  \uses{def:pet-ch8-killing-field,lem:pet-ch8-grove-searle-codim2,lem:pet-ch8-isolated-zero-orbit-angle-bound,lem:pet-ch8-toponogov-angle-lower-bound}
  If $M^4$ is a compact, orientable, positively curved $4$-manifold that admits
  a Killing field, then $\chi(M)\le3$. In particular, $M$ is topologically
  equivalent to $S^4$ or $\mathbb{CP}^2$.
  \source{petersen:page-0336}
\end{theorem}
\begin{proof}
  \uses{thm:pet-ch8-hsiang-kleiner-4manifold}
  We may normalize $\sec\ge1$, so that Toponogov's theorem applies with the
  comparison space $S^2(1)$. If the zero set of the Killing field has a
  $2$-dimensional component, the result follows from
  \cref{lem:pet-ch8-grove-searle-codim2} (applied with $n=4$). Otherwise all
  zeros of $X$ are isolated, and it suffices to derive a contradiction from the
  existence of four (or more) isolated zeros $p_1,p_2,p_3,p_4$.

  By \cref{lem:pet-ch8-isolated-zero-orbit-angle-bound}, at each isolated zero
  $p_i$ the flow of $X$ maps segments between any two other zeros to segments
  between the same two zeros, and for each choice of three of the remaining
  points $q,r,s$ (playing the role of $v_1,v_2,v_3\in S^3\subset T_{p_i}M$ via
  the exponential map at $p_i$) one gets, for any $\varepsilon>0$,
  \[
    \inf\angle q\,p_i\,r+\inf\angle r\,p_i\,s+\inf\angle s\,p_i\,r\le\pi
  \]
  after letting $\varepsilon\to0$ (using density of the orbit as in
  \cref{lem:pet-ch8-isolated-zero-orbit-angle-bound}). Summing this bound over
  all four choices of which point plays the role of $p_i$ (with the other three
  as $q,r,s$) gives a total sum $\le4\pi$.

  On the other hand, \cref{lem:pet-ch8-toponogov-angle-lower-bound} gives, for
  each of the four choices of vertex point, a strict lower bound
  $\inf\angle p\,x\,q+\inf\angle q\,x\,p+\inf\angle p\,q\,x>\pi$ for the
  corresponding triangle's three vertex angles (each such angle bounded below
  by the comparison angle in $S^2(1)$); summing over all four triangles gives a
  total sum $>4\pi$.

  These two bounds on the same quadruple sum of infima of angles contradict
  each other. Hence $X$ can have at most three isolated zeros, so
  $\chi(M)=\sum_i\chi(\{p_i\})\le3$ by \cref{thm:pet-ch8-euler-characteristic-zero-set-sum}
  in the all-isolated case; combined with the codimension-$2$ case handled by
  \cref{lem:pet-ch8-grove-searle-codim2}, this yields $\chi(M)\le3$ in general,
  and by the classification in \cref{lem:pet-ch8-grove-searle-codim2}, $M$ is
  topologically $S^4$ or $\mathbb{CP}^2$.
\end{proof}

\begin{theorem}[Theorem 8.3.5 (Grove and Searle, 1994)]
  \label{thm:pet-ch8-grove-searle-commuting-dependent}
  \lean{PetersenLib.groveSearle_commutingKillingFields_dependentSomewhere}
  \uses{def:pet-ch8-killing-field,prop:pet-ch8-killing-curvature-identities}
  Let $M$ be a compact $n$-manifold with positive sectional curvature. Two
  commuting Killing fields must be linearly dependent somewhere on $M$.
  \source{petersen:page-0337,petersen:page-0338}
\end{theorem}
\begin{proof}
  \uses{thm:pet-ch8-grove-searle-commuting-dependent}
  Let $X,Y$ be commuting Killing fields on $M$. By
  \cref{prop:pet-ch8-killing-curvature-identities}, $[\nabla X,\nabla Y]=R(X,Y)$
  (as $[X,Y]=0$), so
  \begin{align*}
    R(X,Y,Y,X)&=g([\nabla X,\nabla Y](Y),X)\\
    &=g(\nabla_{\nabla_YY}X,X)-g(\nabla_{\nabla_YX}Y,X)\\
    &=-g(\nabla_XX,\nabla_YY)+|\nabla_XY|^2,
  \end{align*}
  where the last step uses skew-symmetry of $\nabla X,\nabla Y$ and
  $g(\nabla_{\nabla_YY}X,X)=-g(\nabla_YY,\nabla_XX)$,
  $g(\nabla_{\nabla_YX}Y,X)=-g(\nabla_YX,\nabla_YX)$ (using also
  $\nabla_YX=\nabla_XY$ since $[X,Y]=0$).

  If $Y$ vanishes somewhere on $M$, the theorem holds. Otherwise consider
  \[
    f=\frac12\Big(|X|^2-\frac{g(X,Y)^2}{|Y|^2}\Big).
  \]
  If $f$ vanishes somewhere, $X$ and $Y$ are linearly dependent there
  (equality in Cauchy--Schwarz), and the theorem holds. Otherwise $f$ attains a
  positive minimum at some $p\in M$. Rescale $Y$ to be a unit vector at $p$ and
  replace $X$ by $\bar X=X-g(X|_p,Y|_p)Y$, so that $\bar X|_p\perp Y|_p$;
  neither change alters the value of $f$ (as $f$ is invariant under adding a
  multiple of $Y$ to $X$ and under rescaling $Y$). Then
  \[
    f=\frac12\Big(|\bar X|^2-\frac{g(\bar X,Y)^2}{|Y|^2}\Big)\le\frac12|\bar X|^2,
  \]
  with equality at $p$; so $\tfrac12|\bar X|^2$ also has a positive minimum at
  $p$, giving $\nabla_{\bar X}\bar X=0$ at $p$ (as in
  \cref{prop:pet-ch8-killing-norm-formulas}(1)). Since $g(\bar X,Y)=0$ at $p$,
  the Hessian of $f$ at $p$ reduces to
  \[
    \mathrm{Hess}\,f(v,v)=|\nabla_v\bar X|^2-R(v,\bar X,\bar X,v)-\big(D_vg(\bar X,Y)\big)^2,
    \qquad v\in T_pM.
  \]
  Using $\nabla Y=\nabla X$ (as $X,Y$ commute) and skew-symmetry,
  \[
    D_vg(\bar X,Y)=g(\nabla_v\bar X,Y)+g(\bar X,\nabla_vY)
    =g(\nabla_v\bar X,Y)-g(v,\nabla_{\bar X}Y)=g(\nabla_v\bar X,2Y).
  \]
  If there is $v\perp\bar X$ at $p$ with $\nabla_v\bar X=0$, then
  $\mathrm{Hess}\,f(v,v)=-R(v,\bar X,\bar X,v)<0$, a contradiction. Otherwise
  $\ker(\nabla\bar X)|_p=\mathrm{span}\{\bar X|_p\}$ (since $\nabla_{\bar X}\bar X=0$
  at $p$), and as $Y\perp\bar X$ at $p$ there is $v\perp\bar X$ with
  $\nabla_v\bar X=2Y$; then $D_vg(\bar X,Y)=4$ and
  \[
    \mathrm{Hess}\,f(v,v)=|2Y|^2-R(v,\bar X,\bar X,v)-16.
  \]
  Choosing $Y$ to have unit length (as arranged) makes $|2Y|^2-16$ appropriately
  bounded so that, together with the strict positivity of
  $R(v,\bar X,\bar X,v)$, $\mathrm{Hess}\,f(v,v)<0$, again contradicting
  $\mathrm{Hess}\,f|_p\ge0$. In either case we obtain a contradiction, so $f$
  must vanish somewhere, and $X,Y$ are linearly dependent there.
\end{proof}

\begin{lemma}[orthogonal fields along a geodesic with drift towards the action]
  \label{lem:pet-ch8-jacobi-type-field}
  \lean{PetersenLib.wilkingOrthogonalField}
  \uses{def:pet-ch8-killing-field}
  Let $M^n$ have positive sectional curvature, $X$ a Killing field, and let
  $N^{n-k}$ be a component of the zero set of $X$. Let $c:[0,1]\to M$ be a unit
  speed geodesic perpendicular to $N$ at both endpoints. Consider vector fields
  $E$ along $c$, orthogonal to both $c$ and the flow direction of $X$, defined
  by the linear ODE
  \[
    E(0)\in T_{c(0)}N,\qquad
    \dot E=-\frac{g(E,\nabla_{\dot c}X)}{|X|^2}X.
  \]
  Then $g(E,X)\equiv0$; $g(E(1),\nabla_{\dot c(1)}X)=0$; $g(E,\dot c)\equiv0$;
  and the space of such fields $E$ for which $E(1)\in T_{c(1)}N$ has dimension
  at least $n-2k+2$.
  \source{petersen:page-0338,petersen:page-0339}
\end{lemma}
\begin{proof}
  \uses{lem:pet-ch8-jacobi-type-field}
  The ODE is non-singular at $t=0$: since
  $g(E,\nabla_{\dot c}X)=-g(\dot c,\nabla_EX)$ (by skew-symmetry of $\nabla X$)
  and $E(0)\in T_{c(0)}N$, evaluating at $t=0$ gives
  $g(E(0),\nabla_{\dot c(0)}X)=-g(\dot c(0),\nabla_{E(0)}X)=0$, since
  $\nabla_{E(0)}X=0$ for $E(0)$ tangent to $N$ (by \cref{thm:pet-ch8-zero-set-totally-geodesic},
  $\nabla X$ restricted to $TN$ vanishes, as $TN=\ker(\nabla X)$ there).

  For $g(E,X)\equiv0$: differentiate,
  \[
    \frac{d}{dt}g(E,X)=g(\dot E,X)+g(E,\nabla_{\dot c}X)
    =-\frac{g(E,\nabla_{\dot c}X)}{|X|^2}g(X,X)+g(E,\nabla_{\dot c}X)=0,
  \]
  and $g(E(0),X(0))=0$ since $X|_{c(0)}=0$ ($c(0)\in N$); hence $g(E,X)\equiv0$
  on $[0,1]$. Since $X|_{c(1)}=0$ as well ($c(1)\in N$), and
  $g(E(1),\nabla_{\dot c(1)}X)=-g(\dot c(1),\nabla_{E(1)}X)$ is finite while the
  differential equation forces $\dot E$ to be a multiple of $X$ (which vanishes
  at $t=1$), a continuity argument gives $g(E(1),\nabla_{\dot c(1)}X)=0$.

  For $g(E,\dot c)\equiv0$: first, $\frac{d}{dt}g(X,\dot c)=g(\nabla_{\dot c}X,\dot c)=0$
  by skew-symmetry (setting $v=w=\dot c$ in
  $g(\nabla_vX,w)=-g(\nabla_wX,v)$), and $g(X,\dot c)|_{t=0}=0$ since
  $X|_{c(0)}=0$; hence $g(X,\dot c)\equiv0$. Then
  \[
    \frac{d}{dt}g(E,\dot c)=g(\dot E,\dot c)=-\frac{g(E,\nabla_{\dot c}X)}{|X|^2}g(X,\dot c)=0,
  \]
  and $g(E(0),\dot c(0))=0$ since $\dot c(0)\perp T_{c(0)}N$ and $E(0)\in T_{c(0)}N$;
  hence $g(E,\dot c)\equiv0$.

  Finally, $E(1)\perp\mathrm{span}\{\dot c(1),\nabla_{\dot c(1)}X\}$, and this
  span is $2$-dimensional (as $\dot c(1)\perp T_{c(1)}N$ so $\dot c(1)\notin\ker(\nabla X)|_{c(1)}$,
  and $\nabla_{\dot c(1)}X\ne0$, $\nabla_{\dot c(1)}X\perp\dot c(1)$ by
  skew-symmetry). The linear map $E(0)\mapsto E(1)$ from $T_{c(0)}N$
  ($(n-k)$-dimensional) into the orthogonal complement of $\mathrm{span}\{\dot c(1),\nabla_{\dot c(1)}X\}$
  in the full $n$-dimensional space has image of dimension at least
  $(n-k)-2$ once further restricted to land inside the $(n-k)$-dimensional
  $T_{c(1)}N$; more precisely, requiring $E(1)\in T_{c(1)}N$ (which has
  codimension $k$, further intersected with the already-forced $2$-codimensional
  orthogonality constraint) leaves a space of solutions $E$ of dimension at
  least $(n-k)-2(k-1)=n-2k+2$.
\end{proof}

\begin{lemma}[Lemma 8.3.6 (connectedness principle with symmetries, Wilking, 2003)]
  \label{lem:pet-ch8-wilking-connectedness-principle}
  \lean{PetersenLib.wilkingConnectednessPrincipleWithSymmetries}
  \uses{def:pet-ch8-killing-field,lem:pet-ch8-jacobi-type-field,thm:pet-ch8-zero-set-totally-geodesic}
  Let $M^n$ be a compact $n$-manifold with positive sectional curvature and $X$
  a Killing field. If $N^{n-k}$ is a component of the zero set of $X$, then
  $N\subset M$ is $(n-2k+2)$-connected.
  \source{petersen:page-0338,petersen:page-0339,petersen:page-0340}
\end{lemma}
\begin{proof}
  \uses{lem:pet-ch8-wilking-connectedness-principle}
  Let $c$ be a unit-speed geodesic perpendicular to $N$ at both endpoints. By
  \cref{lem:pet-ch8-jacobi-type-field} there is a space of dimension at least
  $n-2k+2$ of fields $E$ along $c$, orthogonal to $\dot c$ and to $X$, with
  $E(0),E(1)\in TN$ and satisfying the linear ODE
  $\dot E=-\tfrac{g(E,\nabla_{\dot c}X)}{|X|^2}X$.

  For $\lambda>0$, let $g_\lambda$ be the Cheeger-deformation-type family of
  metrics on $M$ obtained by squeezing the size of $X$ towards $0$ as
  $\lambda\to0$, while leaving directions orthogonal to $X$ unchanged; $c$
  remains a $g_\lambda$-geodesic since $\nabla_YY$ for $Y\perp X$ is unaffected
  by this deformation. The second variation of energy for such $E$, in the
  metric $g_\lambda$, is
  \begin{align*}
    \left.\frac{d^2E_\lambda}{ds^2}\right|_{s=0}
    &=\int_0^1|\dot E|_{g_\lambda}^2\,dt-\int_0^1g_\lambda(R(E,\dot c)\dot c,E)\,dt\\
    &\le\int_0^1\Big|\frac{g(E,\nabla_{\dot c}X)}{|X|_\lambda^2}X\Big|^2\,dt
      -\int_0^1\sec_g(E,\dot c)|E|_\lambda^2\,dt\\
    &=\int_0^1\frac{|g(E,\nabla_{\dot c}X)|^2}{|X|_g^4}|X|_{g_\lambda}^2\,dt
      -\int_0^1\sec_g(E,\dot c)|E|_\lambda^2\,dt\\
    &\longrightarrow-\int_0^1\sec_g(E,\dot c)|E|_\lambda^2\,dt\quad\text{as }\lambda\to0,
  \end{align*}
  since $|X|_{g_\lambda}\to0$ while $g(E,\nabla_{\dot c}X)/|X|_g^2$ stays
  bounded. As $M$ has positive sectional curvature and $E,\dot c$ are
  orthogonal to $X$ (hence their sectional curvature $\sec_g(E,\dot c)$ is
  unaffected by the deformation and remains positive, in fact $\sec_{g_\lambda}(E,\dot c)$
  can only increase as $\lambda\to0$), the limiting second variation is
  strictly negative for every nonzero $E$ in this space, for all sufficiently
  small $\lambda$. Since this holds simultaneously for the whole
  $(n-2k+2)$-dimensional space of such fields, and independently of the chosen
  $c$ perpendicular to $N$ at its endpoints, in the metric $g_\lambda$ (for
  $\lambda$ small) every such geodesic $c$ has Morse index at least $n-2k+2$.
  By the standard relationship between geodesic index bounds relative to a
  submanifold and homotopy connectivity, this shows $N\subset M$ is
  $(n-2k+2)$-connected.
\end{proof}

\begin{lemma}[Lemma 8.3.7 (Grove and Searle, 1994)]
  \label{lem:pet-ch8-grove-searle-codim2}
  \lean{PetersenLib.groveSearle_codimTwoZeroSet_classification}
  \uses{lem:pet-ch8-wilking-connectedness-principle,def:pet-ch8-euler-characteristic}
  Let $M$ be a closed $n$-manifold with positive sectional curvature. If $M$
  admits a Killing field whose zero set has a component $N$ of codimension
  $2$, then $M$ is diffeomorphic to $S^n$, $\mathbb{CP}^{n/2}$, or a cyclic
  quotient of a sphere $S^n/\mathbb{Z}_q$.
  \source{petersen:page-0340,petersen:page-0341}
\end{lemma}
\begin{proof}
  \uses{lem:pet-ch8-grove-searle-codim2}
  We prove only the (co)homology statement for simply connected $M$. By
  \cref{lem:pet-ch8-wilking-connectedness-principle} with $k=2$, $N\subset M$ is
  $(n-2)$-connected, so for $k<n-2$ the maps $H_k(N)\to H_k(M)$ and
  $H^k(M)\to H^k(N)$ are isomorphisms. Combining this with Poincar\'e duality
  $H_k(M)\simeq H^{n-k}(M)$ and $H_k(N)\simeq H^{n-2-k}(N)$ gives, for
  $0<k<n-2$, chains of isomorphisms
  \[
    H_{k+2}(M)\to H^{n-k-2}(M)\to H^{n-k-2}(N)\to H_k(N)\to H_k(M),
  \]
  \[
    H^k(M)\to H^k(N)\to H_{n-2-k}(N)\to H_{n-2-k}(M)\to H^{k+2}(M).
  \]
  Since $M$ is simply connected, $H_1(M)=0$. Using the first chain of
  isomorphisms iteratively: if $n$ is even, one obtains
  $0\simeq H_1(M)\simeq H_3(M)\simeq\cdots\simeq H_{n-1}(M)$; if $n$ is odd,
  \[
    0\simeq H_1(M)\simeq H_3(M)\simeq\cdots\simeq H_{n-2}(M)\simeq H^2(M)\simeq\cdots\simeq H^{n-1}(M).
  \]
  In particular, when $n$ is odd, all of $H_1(M),\dots,H_{n-1}(M)$ vanish (using
  Poincar\'e duality to also fold in the even-degree groups), so $M$ has the
  integral homology of $S^n$ — a homotopy sphere.

  When $n$ is even, the second chain of isomorphisms, combined with
  $H^{n-2}(N)\simeq\mathbb{Z}$ (as $N$ has codimension $2$, so is
  even-dimensional of dimension $n-2$, and is itself simply connected and
  positively curved, hence orientable) and injectivity of
  $H^{n-2}(M)\to H^{n-2}(N)$ (from $(n-2)$-connectivity of $N\subset M$), gives
  \[
    H^2(M)\simeq H^4(M)\simeq\cdots\simeq H^{n-2}(M)\hookrightarrow H^{n-2}(N)\simeq\mathbb{Z}.
  \]
  Hence these even-degree cohomology groups are all trivial or all isomorphic
  to $\mathbb{Z}$: in the first case $M$ has the cohomology of $S^n$, in the
  second that of $\mathbb{CP}^{n/2}$. (When $M$ is not simply connected the
  same argument applies to the universal cover, identifying $M$ with a
  quotient, and further work isolates the possible cyclic quotients
  $S^n/\mathbb{Z}_q$.)
\end{proof}

\begin{proposition}[Proposition 8.3.8]
  \label{prop:pet-ch8-commuting-killing-fields-properties}
  \lean{PetersenLib.commutingKillingFields_tangentToLevelSets,PetersenLib.commutingKillingFields_linearCombinationVanishesOnLarger}
  \uses{def:pet-ch8-killing-field,prop:pet-ch8-killing-curvature-identities}
  Let $(M,g)$ be compact and $X,Y\in\iso(M,g)$ commute.
  \begin{enumerate}
    \item $Y$ is tangent to the level sets of $|X|^2$, in particular to the
      zero set of $X$.
    \item If $X$ and $Y$ both vanish on a totally geodesic submanifold
      $N\subset M$, then some linear combination of $X,Y$ vanishes on a
      strictly larger connected submanifold.
  \end{enumerate}
  \source{petersen:page-0341}
\end{proposition}
\begin{proof}
  \uses{prop:pet-ch8-commuting-killing-fields-properties}
  (1) Since $[X,Y]=0$, $L_YX=0$, so
  \[
    0=(L_Yg)(X,X)=D_Y|X|^2-2g(L_YX,X)=D_Y|X|^2,
  \]
  using that $L_Yg=0$ as $Y$ is a Killing field. Hence the flow of $Y$
  preserves the level sets of $|X|^2$, in particular the zero set (level
  $0$).

  (2) We may take $N$ to have even codimension (by
  \cref{thm:pet-ch8-zero-set-totally-geodesic}). Fix $p\in N$; by
  \cref{prop:pet-ch8-killing-curvature-identities}, $[\nabla X,\nabla Y]=R(X,Y)$
  (since $[X,Y]=0$), and evaluated at $p$ where $X|_p=Y|_p=0$ this gives
  $[(\nabla X)|_p,(\nabla Y)|_p]=0$: the two skew-symmetric endomorphisms
  commute. Hence they may be simultaneously block-diagonalized, giving a
  splitting
  \[
    T_pM=T_pN\oplus E_1\oplus\cdots\oplus E_k,
  \]
  with each $E_i$ a $2$-dimensional invariant subspace for both $(\nabla X)|_p$
  and $(\nabla Y)|_p$. Since the space of skew-symmetric endomorphisms of a
  $2$-dimensional space is $1$-dimensional, on $E_1$ there is a linear
  combination $\alpha(\nabla X)|_p+\beta(\nabla Y)|_p$ (not both $\alpha,\beta$
  zero) that vanishes on $E_1$. Let $\bar N\supset N$ be the connected
  component, containing $p$, of the zero set of $\alpha X+\beta Y$. Then
  $T_p\bar N=\ker(\nabla(\alpha X+\beta Y))|_p\supset T_pN\oplus E_1$
  (strictly larger than $T_pN$), so $N\subsetneq\bar N$.
\end{proof}

\begin{definition}[symmetry rank]
  \label{def:pet-ch8-symmetry-rank}
  \lean{PetersenLib.symmetryRank}
  \uses{def:pet-ch8-killing-field}
  The \emph{symmetry rank} of a Riemannian manifold $(M,g)$ is the dimension of
  a maximal subspace (equivalently, Abelian subalgebra) of commuting Killing
  fields in $\iso(M,g)$.
  \source{petersen:page-0341}
\end{definition}

\begin{theorem}[Theorem 8.3.9 (Grove and Searle, 1994)]
  \label{thm:pet-ch8-symmetry-rank-half-classification}
  \lean{PetersenLib.groveSearle_symmetryRank_atLeastHalf_classification}
  \uses{def:pet-ch8-symmetry-rank,lem:pet-ch8-grove-searle-codim2,prop:pet-ch8-commuting-killing-fields-properties}
  Let $M$ be a compact $n$-manifold with positive sectional curvature and
  symmetry rank $k$. If $k\ge n/2$, then $M$ is diffeomorphic to a sphere, a
  complex projective space, or a cyclic quotient of a sphere $S^n/\mathbb{Z}_q$,
  with $\mathbb{Z}_q$ acting by isometries on the unit sphere.
  \source{petersen:page-0341,petersen:page-0342}
\end{theorem}
\begin{proof}
  \uses{thm:pet-ch8-symmetry-rank-half-classification}
  Choose an Abelian subalgebra $\mathfrak a\subset\iso(M,g)$ of dimension
  $k\ge n/2$. \cref{prop:pet-ch8-commuting-killing-fields-properties}(2)
  (applied repeatedly, starting from the zero set of a single $X\in\mathfrak a$
  and enlarging by successive linear combinations of commuting fields in
  $\mathfrak a$) produces a maximal (with respect to inclusion), nontrivial,
  totally geodesic submanifold $N\subset M$ that is the zero set of some
  $X\in\mathfrak a$, with $\dim(\mathfrak a|_N)=k-1$, where $\mathfrak a|_N$
  denotes the restriction of the fields in $\mathfrak a$ tangent to $N$
  (\cref{prop:pet-ch8-commuting-killing-fields-properties}(1) shows every
  $Y\in\mathfrak a$ is tangent to $N$).

  If $N$ does not have codimension $2$, this construction may be iterated
  inside $N$ to produce a totally geodesic $1$- or $2$-dimensional submanifold
  $S\subset M$ with $\dim(\mathfrak a|_S)\ge2$. A $1$-manifold has
  $1$-dimensional isometry group, so this case does not occur. If $\dim S=2$,
  pick $X,Y\in\mathfrak a|_S$; since $X$ is nontrivial on the compact surface
  $S$ it has an isolated zero $p$ (a positively curved surface's Killing fields
  vanish at isolated points), and $Y$, preserving the (single-point) zero-set
  component $\{p\}$ of $X$ by
  \cref{prop:pet-ch8-commuting-killing-fields-properties}(1), also vanishes at
  $p$. Then $(\nabla X)|_p,(\nabla Y)|_p:T_pM\to T_pM$ determine $X,Y$
  completely (by \cref{prop:pet-ch8-killing-uniqueness}), and since the space
  of skew-symmetric endomorphisms of $T_pM$ (a $2$-plane if $S=T_pM$ locally
  agrees with the ambient tangent directions relevant here) is $1$-dimensional,
  $X,Y$ are linearly dependent, forcing $\dim(\mathfrak a|_S)=1$, a
  contradiction. Hence $N$ has codimension $2$, and
  \cref{lem:pet-ch8-grove-searle-codim2} finishes the proof.
\end{proof}

\begin{theorem}[Theorem 8.3.10 (Puttmann and Searle, 2002 and Rong, 2002)]
  \label{thm:pet-ch8-puttmann-searle-rong}
  \lean{PetersenLib.puttmannSearleRong_symmetryRank_eulerChar_pos}
  \uses{def:pet-ch8-symmetry-rank,lem:pet-ch8-grove-searle-codim2,cor:pet-ch8-positive-euler-characteristic-6manifold}
  If $M^{2n}$ is a compact $2n$-manifold with positive sectional curvature and
  symmetry rank $k\ge n/2-1$, then $\chi(M)>0$.
  \source{petersen:page-0342}
\end{theorem}
\begin{proof}
  \uses{thm:pet-ch8-puttmann-searle-rong}
  For $2n=2,4$ there is no constraint on $k$ and the statement is
  \cref{rem:pet-ch8-hopf-conjecture-low-dimensions}; for $2n=6$ it is
  \cref{cor:pet-ch8-positive-euler-characteristic-6manifold} (Berger's
  theorem). For $2n=8$ the argument is as in the $6$-dimensional case unless
  the zero set of some Killing field has a $6$-dimensional component, in which
  case \cref{lem:pet-ch8-grove-searle-codim2} (applied to that component, whose
  codimension in $M$ is $2$) establishes the claim.

  In general, induction on $n$ requires the stronger statement: if
  $\mathfrak a\subset\iso(M,g)$ is Abelian of dimension $k\ge n/2-1$, then every
  component of the zero set of any $X\in\mathfrak a$ has positive Euler
  characteristic. This holds for $2n=2,4,6,8$ as above.

  For the induction step there are two cases. If every zero set $N\subset M$ of
  every $X\in\mathfrak a$ has codimension $\ge4$: taking $N$ maximal and
  nontrivial gives (as in the proof of
  \cref{thm:pet-ch8-symmetry-rank-half-classification}) $\dim(\mathfrak a|_N)\ge k-1$
  and $\dim N\ge2(k-1)$, so $N$ satisfies the stronger induction hypothesis in
  its own (lower) dimension, and since every component of a zero set is
  contained in some maximal nontrivial such $N$, the induction hypothesis
  applies to conclude that component has positive Euler characteristic.

  If instead some zero set $N$ for $X\in\mathfrak a$ has codimension $2$, then
  \cref{lem:pet-ch8-grove-searle-codim2} shows the odd Betti numbers of $M$
  vanish, which in turn (by the same connectedness and duality arguments)
  forces the odd Betti numbers of every component of every zero set of every
  Killing field on $M$ to vanish, giving positive Euler characteristic for
  those components as well.
\end{proof}

\begin{remark}[a positively curved exception with lower symmetry rank]
  \label{rem:pet-ch8-f4-spin8-exception}
  \lean{PetersenLib.f4Spin8SymmetryRankException}
  \uses{def:pet-ch8-symmetry-rank}
  \cref{thm:pet-ch8-puttmann-searle-rong} does not cover all known examples:
  the homogeneous space $F_4/\mathrm{Spin}(8)$ is a positively curved
  $24$-manifold with symmetry rank $4$.
  \source{petersen:page-0342}
\end{remark}

\begin{theorem}[Theorem 8.3.11 (Wilking, 2003)]
  \label{thm:pet-ch8-wilking-symmetry-rank-classification}
  \lean{PetersenLib.wilking_symmetryRank_classification}
  \uses{def:pet-ch8-symmetry-rank}
  Let $M$ be a compact, simply connected, positively curved $n$-manifold with
  symmetry rank $k$, and $n\ge10$. If $k\ge n/4+1$, then $M$ has the topology
  of a sphere, complex projective space, or quaternionic projective space.
  Moreover, when $M$ is not simply connected, its fundamental group is cyclic.
  \source{petersen:page-0343}
\end{theorem}

\begin{theorem}[Theorem 8.3.12 (Kennard, 2013)]
  \label{thm:pet-ch8-kennard-2013-euler-characteristic}
  \lean{PetersenLib.kennard2013_symmetryRank_eulerChar_pos}
  \uses{def:pet-ch8-symmetry-rank}
  If $M^{4n}$ is a compact $4n$-manifold with positive sectional curvature and
  symmetry rank $k>2\log_4(4n)-2$, then $\chi(M)>0$.
  \source{petersen:page-0343}
\end{theorem}

\begin{theorem}[Theorem 8.3.13 (Kennard, 2012 and Amann--Kennard, 2014)]
  \label{thm:pet-ch8-kennard-amann-betti-numbers}
  \lean{PetersenLib.kennardAmannKennard_bettiNumberBounds}
  \uses{def:pet-ch8-symmetry-rank}
  Let $M^{2n}$ be a compact positively curved manifold with symmetry rank
  $k\ge\log_{4/3}(2n)$.
  \begin{enumerate}
    \item The Betti numbers satisfy $b_2\le1$ and $b_4\le1$.
    \item If $b_4=0$, then $\chi(M)=2$.
    \item Consequently, $M$ cannot have the homotopy type of a product
      $N\times N$ for compact simply connected $N$.
  \end{enumerate}
  \source{petersen:page-0343}
\end{theorem}

\section{Exercises}

\begin{remark}[Exercise 8.4.1]
  \label{rem:pet-ch8-ex-8-4-1}
  \lean{PetersenLib.exercise_8_4_1}
  \uses{thm:pet-ch8-isometry-lie-algebra-dimension}
  Show that \cref{thm:pet-ch8-isometry-lie-algebra-dimension} does not
  necessarily extend to incomplete Riemannian manifolds.
  \source{petersen:page-0343}
\end{remark}

\begin{remark}[Exercise 8.4.2]
  \label{rem:pet-ch8-ex-8-4-2}
  \lean{PetersenLib.exercise_8_4_2}
  \uses{def:pet-ch8-killing-field,thm:pet-ch8-zero-set-totally-geodesic}
  Let $N$ be a component of the zero set of a Killing field $X$. Show that
  $\nabla_V(\nabla X)=0$ for vector fields $V$ tangent to $N$.
  \source{petersen:page-0343}
\end{remark}

\begin{remark}[Exercise 8.4.3]
  \label{rem:pet-ch8-ex-8-4-3}
  \lean{PetersenLib.exercise_8_4_3}
  \uses{def:pet-ch8-killing-field}
  Show that a coordinate vector field $\partial_i$ is a Killing field if and
  only if $\partial_kg_{ij}=0$.
  \source{petersen:page-0343}
\end{remark}

\begin{remark}[Exercise 8.4.4]
  \label{rem:pet-ch8-ex-8-4-4}
  \lean{PetersenLib.exercise_8_4_4}
  \uses{def:pet-ch8-killing-field}
  Let $X$ be a Killing field on $(M,g)$ and $N\subset M$ a submanifold with the
  induced metric.
  \begin{enumerate}
    \item Show that if $X$ is tangent to $N$, then $X|_N$ is a Killing field
      on $N$.
    \item Show that if $N$ is totally geodesic, then $X^\top$, the projection
      of $X$ onto $N$, is a Killing field.
  \end{enumerate}
  \source{petersen:page-0343}
\end{remark}

\begin{remark}[Exercise 8.4.5]
  \label{rem:pet-ch8-ex-8-4-5}
  \lean{PetersenLib.exercise_8_4_5}
  \uses{def:pet-ch8-killing-field}
  Let $(M,g)$ be a complete Riemannian manifold and $X$ a Killing field on $M$.
  \begin{enumerate}
    \item Let $c:[a,b]\to M$ be a geodesic. Show that there is $\varepsilon>0$
      and a geodesic variation $\bar c:(-\varepsilon,\varepsilon)\times[a,b]\to M$
      such that the curves $s\mapsto\bar c(s,t)$ are integral curves of $X$.
    \item Use completeness to extend $\bar c$ to
      $(-\varepsilon,\varepsilon)\times\mathbb{R}\to M$, and show that for all
      $t\in\mathbb{R}$ the curves $s\mapsto\bar c(s,t)$ are integral curves of
      $X$.
    \item Show that the integral curve of $X$ through any point of $M$ exists
      on $(-\varepsilon,\varepsilon)$.
    \item Show that $X$ is complete.
  \end{enumerate}
  \source{petersen:page-0343,petersen:page-0344}
\end{remark}

\begin{remark}[Exercise 8.4.6]
  \label{rem:pet-ch8-ex-8-4-6}
  \lean{PetersenLib.exercise_8_4_6}
  \uses{def:pet-ch8-symmetry-rank,thm:pet-ch8-bochner-negative-ricci}
  Let $(M^n,g)$ be a compact Riemannian $n$-manifold with $\dim(\iso(M,g))=n$
  and $\mathrm{Ric}\le0$. Show that $M$ is flat and $\pi_1=\mathbb{Z}^n$. (Hint:
  show $\widetilde M=\mathbb{R}^n$ and that the deck transformations have
  differential $I$, since they preserve a parallel orthonormal frame.)
  \source{petersen:page-0344}
\end{remark}

\begin{remark}[Exercise 8.4.7]
  \label{rem:pet-ch8-ex-8-4-7}
  \lean{PetersenLib.exercise_8_4_7}
  \uses{def:pet-ch8-killing-field}
  Let $x^i$ be standard Cartesian coordinates on $\mathbb{R}^n$, and let
  $W=\mathrm{span}_{\mathbb{R}}\{1,x^1,\dots,x^n\}$,
  $V=\mathrm{span}_{\mathbb{R}}\{x^1,\dots,x^n\}$.
  \begin{enumerate}
    \item Show that $\iso(\mathbb{R}^n)$ is naturally isomorphic to $\Lambda^2W$,
      identifying $u\wedge v$ with the vector field $u\nabla v-v\nabla u$.
    \item Show similarly that $\iso(S^{n-1}(R))$ is naturally isomorphic to
      $\Lambda^2V$.
    \item Using $\nabla u=g^{ij}\partial_iu\,\partial_j$ for a
      pseudo-Riemannian space, redo (1) for $\mathbb{R}^{p,q}$ and (2) for
      $H^{n-1}(R)\subset\mathbb{R}^{n,1}$.
  \end{enumerate}
  \source{petersen:page-0344}
\end{remark}

\begin{remark}[Exercise 8.4.8]
  \label{rem:pet-ch8-ex-8-4-8}
  \lean{PetersenLib.exercise_8_4_8}
  \uses{def:pet-ch8-killing-field}
  Let $x^i$ be standard Cartesian coordinates on $\mathbb{R}^{n+1}$ and
  $V=\mathrm{span}_{\mathbb{R}}\{1,x^1,\dots,x^{n+1}\}$, restricted to $S^n$.
  Show that for $u,v\in V$ the fields $u\nabla v-v\nabla u$ are conformal (a
  field $X$ is conformal if $L_Xg=\lambda g$ for some function $\lambda$).
  \source{petersen:page-0344}
\end{remark}

\begin{remark}[Exercise 8.4.9]
  \label{rem:pet-ch8-ex-8-4-9}
  \lean{PetersenLib.exercise_8_4_9}
  \uses{prop:pet-ch8-killing-uniqueness}
  Let $(M,g)$ be a Riemannian manifold and $t_p=\{X|_p\in T_pM\mid X\in\iso\}$.
  \begin{enumerate}
    \item Show that $t_p$ defines an integrable distribution on an open set
      $O\subset M$. (Hint: show
      $\{p\in M\mid\dim t_p=\max_{q\in M}\dim t_q\}$ is open.)
    \item Give an example where $O\ne M$.
    \item Show that $DF(t_p)=t_{F(p)}$ for all $F\in\mathrm{Iso}(M,g)$.
    \item Assume $(M,g)$ is complete. Show that the leaves of this
      distribution are homogeneous and properly embedded. (Hint: show the
      connected subgroup $\mathrm{Iso}_0\subset\mathrm{Iso}(M,g)$ containing
      the identity is closed and preserves the leaves.)
    \item Show that if $(M,g)$ is homogeneous, then $t_p=T_pM$ for all
      $p\in M$.
  \end{enumerate}
  \source{petersen:page-0344}
\end{remark}

\begin{remark}[Exercise 8.4.10]
  \label{rem:pet-ch8-ex-8-4-10}
  \lean{PetersenLib.exercise_8_4_10}
  \uses{def:pet-ch8-killing-field}
  Let $\mathfrak t\subset\iso(M,g)$ be an Abelian subalgebra corresponding to a
  torus subgroup $T^k\subset\mathrm{Iso}(M,g)$. Define $\mathfrak p\subset\mathfrak t$
  as the Killing fields corresponding to circle actions, i.e.\ actions induced
  by homomorphisms $S^1\to T^k$. Show that $\mathfrak p$ is a vector space over
  $\mathbb{Q}$ with $\dim_{\mathbb{Q}}\mathfrak p=\dim_{\mathbb{R}}\mathfrak t$.
  \source{petersen:page-0344}
\end{remark}

\begin{remark}[Exercise 8.4.11]
  \label{rem:pet-ch8-ex-8-4-11}
  \lean{PetersenLib.exercise_8_4_11}
  \uses{def:pet-ch8-killing-field}
  Given two Killing fields $X,Y$ on a Riemannian manifold, develop a formula
  for $\Delta g(X,Y)$. Use this to give a formula for the Ricci curvature in a
  frame consisting of Killing fields.
  \source{petersen:page-0344}
\end{remark}

\begin{remark}[Exercise 8.4.12]
  \label{rem:pet-ch8-ex-8-4-12}
  \lean{PetersenLib.exercise_8_4_12}
  \uses{def:pet-ch8-killing-field}
  Let $X$ be a vector field on a Riemannian manifold.
  \begin{enumerate}
    \item Show that
      $|L_Xg|^2=2|\nabla X|^2+2\operatorname{tr}(\nabla X\circ\nabla X)$.
    \item Establish, on a closed oriented Riemannian manifold,
      \[
        \int_M\Big(\mathrm{Ric}(X,X)+\operatorname{tr}(\nabla X\circ\nabla X)-(\operatorname{div}X)^2\Big)=0,
      \]
      \[
        \int_M\Big(\mathrm{Ric}(X,X)+g(\operatorname{tr}\nabla^2X,X)+\tfrac12|L_Xg|^2-(\operatorname{div}X)^2\Big)=0.
      \]
    \item Show that $X$ is a Killing field if and only if
      $\operatorname{div}X=0$ and $\operatorname{tr}\nabla^2X=-\mathrm{Ric}(X)$.
  \end{enumerate}
  \source{petersen:page-0345}
\end{remark}

\begin{remark}[Exercise 8.4.13 (Yano)]
  \label{rem:pet-ch8-ex-8-4-13}
  \lean{PetersenLib.exercise_8_4_13}
  \uses{rem:pet-ch8-ex-8-4-12}
  If $X$ is an affine vector field, show that
  $\operatorname{tr}\nabla^2X=-\mathrm{Ric}(X)$ and that $\operatorname{div}X$
  is constant. Use this together with the characterizations of
  \cref{rem:pet-ch8-ex-8-4-12} to show that on closed manifolds affine fields
  are Killing fields.
  \source{petersen:page-0345}
\end{remark}

\begin{remark}[Exercise 8.4.14]
  \label{rem:pet-ch8-ex-8-4-14}
  \lean{PetersenLib.exercise_8_4_14}
  \uses{def:pet-ch8-killing-field}
  Let $X$ be a vector field on a Riemannian manifold. Show that $X$ is a
  Killing field if and only if $L_X$ and $\Delta$ commute on functions.
  \source{petersen:page-0345}
\end{remark}

\begin{remark}[Exercise 8.4.15]
  \label{rem:pet-ch8-ex-8-4-15}
  \lean{PetersenLib.exercise_8_4_15}
  \uses{def:pet-ch8-symmetry-rank}
  Let $(M,g)$ be a compact $n$-manifold with positive sectional curvature and
  $\mathfrak a\subset\iso$ an Abelian subalgebra.
  \begin{enumerate}
    \item Show that $\dim\mathfrak a\le n/2$.
    \item Show that spheres and complex projective spaces have maximal
      symmetry rank.
    \item Show that the flat torus $T^n$ has symmetry rank $n$.
  \end{enumerate}
  \source{petersen:page-0345}
\end{remark}

\begin{remark}[Exercise 8.4.16]
  \label{rem:pet-ch8-ex-8-4-16}
  \lean{PetersenLib.exercise_8_4_16}
  \uses{def:pet-ch8-symmetry-rank,prop:pet-ch8-commuting-killing-fields-properties}
  Let $(M,g)$ be a compact $n$-manifold with positive sectional curvature and
  $\mathfrak a\subset\iso$ an Abelian subalgebra. Let $\mathcal Z(\mathfrak a)$
  be the set of nontrivial connected components of zero sets of Killing fields
  in $\mathfrak a$. Show that if $N\in\mathcal Z(\mathfrak a)$ is maximal under
  inclusion, then $\dim(\mathfrak a|_N)=\dim\mathfrak a-1$ and
  $\dim N\ge2(\dim\mathfrak a-1)$.
  \source{petersen:page-0345}
\end{remark}

\begin{remark}[Exercise 8.4.17 (Kennard)]
  \label{rem:pet-ch8-ex-8-4-17}
  \lean{PetersenLib.exercise_8_4_17}
  \uses{def:pet-ch8-symmetry-rank,thm:pet-ch8-symmetry-rank-half-classification,rem:pet-ch8-ex-8-4-16,lem:pet-ch8-wilking-connectedness-principle}
  Let $(M^n,\rho)$ be a simply connected compact $n$-manifold with positive
  sectional curvature and symmetry rank $\ge3n/8+1$.
  \begin{enumerate}
    \item When $n\le8$, conclude that the assumptions are covered by
      \cref{thm:pet-ch8-symmetry-rank-half-classification}.
    \item Use \cref{rem:pet-ch8-ex-8-4-16} to show that if
      $N\in\mathcal Z(\mathfrak a)$ is maximal under inclusion, then
      $\dim N\ge3n/4$.
    \item Show that a maximal $N\in\mathcal Z(\mathfrak a)$ either has
      codimension $2$ or has symmetry rank $\ge\tfrac{3\dim N}8+1$.
    \item Use induction on $n$ to show that $M$ has the homology groups of a
      sphere or complex projective space. (Hint: when the maximal
      $N\in\mathcal Z(\mathfrak a)$ has codimension $\ge4$, use
      \cref{lem:pet-ch8-wilking-connectedness-principle} to show $N$ is simply
      connected, then use it again to compute homology/homotopy groups in
      dimensions $\le n/2$; finally use Poincar\'e duality to find all
      homology groups of $M$.)
  \end{enumerate}
  \source{petersen:page-0345,petersen:page-0346}
\end{remark}

\begin{remark}[Exercise 8.4.18]
  \label{rem:pet-ch8-ex-8-4-18}
  \lean{PetersenLib.exercise_8_4_18}
  \uses{def:pet-ch8-killing-field}
  Let $\widetilde M\to M$ be the universal covering, with deck transformations
  $\Gamma=\pi_1(M)$ acting as isometries on $\widetilde M$.
  \begin{enumerate}
    \item Show that
      $\iso(M)=\{X\in\iso(\widetilde M)\mid G_*X=X\text{ for all }G\in\Gamma\}$.
    \item Show that the connected Lie subgroup corresponding to
      $\iso(M)\subset\iso(\widetilde M)$ is the identity component of the
      centralizer
      $\mathbf C(\Gamma)=\{F\in\mathrm{Iso}(\widetilde M)\mid FG=GF\text{ for all }G\in\Gamma\}$.
    \item Show that $\mathrm{Iso}(M)$ can be identified with
      $\mathbf N(\Gamma)/\Gamma$, where
      $\mathbf N(\Gamma)=\{F\in\mathrm{Iso}(\widetilde M)\mid F\Gamma F^{-1}=\Gamma\}$
      is the normalizer.
  \end{enumerate}
  \source{petersen:page-0346}
\end{remark}
